\documentclass[11pt]{article}

% ---------- Packages ----------
\usepackage[margin=1in]{geometry}
\usepackage{amsmath, amssymb, amsthm}
\usepackage{mathtools}

% ---------- Macros ----------
\newcommand{\E}{\mathbb{E}}
\newcommand{\R}{\mathbb{R}}
\newcommand{\cD}{\mathcal{D}}
\newcommand{\cF}{\mathcal{F}}
\newcommand{\cR}{\mathcal{R}}
\newcommand{\argmin}{\operatorname*{arg\,min}}

% ---------- Theorem environments ----------
\newtheorem{theorem}{Theorem}[section]
\newtheorem{proposition}[theorem]{Proposition}

\title{Cross-Validation Under Regime Shifts}
\author{Yegor Baranovski}
\date{\today}

\begin{document}
\maketitle

\section{Problem statement}
Many prediction problems involve data collected across multiple regimes (e.g., time periods or macroeconomic
conditions). In such settings, the joint distribution of $(X,Y)$ may change across regimes, so observations are not
identically distributed over the full sample. Nevertheless, standard cross-validation (especially random $K$-fold CV)
is typically justified under i.i.d./exchangeability-type assumptions. When regimes are mixed across folds, CV can
estimate the wrong generalization target, leading to optimistic performance estimates, unstable hyperparameter
selection, and even qualitative failures such as coefficient sign flips across time.



\section{Deployment objectives under regime shift}
When data arise from multiple regimes, the phrase ``best model'' is ambiguous without specifying a deployment
objective. Three common targets are:
\begin{enumerate}
  \item \textbf{Regime-specific deployment:} choose a model that performs best on a particular regime of interest
  (often a future regime).
  \item \textbf{Robust deployment:} choose a model that performs well across regimes, e.g., by minimizing the worst-case
  risk over regimes.
  \item \textbf{Mixture deployment:} choose a model that performs best under a specified mixture of regimes, reflecting
  anticipated deployment frequencies.
\end{enumerate}
Here, we will focus on the \emph{regime-specific deployment} objective. This choice yields a clean
mathematical target (train on past regimes, deploy on a designated regime) and matches common practice in
applications where models are trained on historical data and deployed in a specific future period.




\section{Setup and goal}

\paragraph{Training and deployment distributions.}
Let $(\mathcal X,\mathcal B_{\mathcal X})$ be the feature space equipped with a $\sigma$-algebra
$\mathcal B_{\mathcal X}$,
and let $(\mathcal Y,\mathcal B_{\mathcal Y})$ be the outcome space with its own $\sigma$-algebra.
The product space $(\mathcal X\times\mathcal Y,\ \mathcal B_{\mathcal X}\otimes\mathcal B_{\mathcal Y})$
represents the joint space of feature--outcome pairs, where
$\mathcal B_{\mathcal X}\otimes\mathcal B_{\mathcal Y}$ is the product $\sigma$-algebra generated by
sets of the form $A\times B$ with $A\in\mathcal B_{\mathcal X}$ and $B\in\mathcal B_{\mathcal Y}$.
A probability measure on this space assigns a probability to each measurable set of
feature--outcome pairs; crucially, the same underlying space can support multiple probability
measures, each reflecting a different data-generating process.
For each regime $r$, let $P_r$ be such a probability measure on
$(\mathcal X\times\mathcal Y,\ \mathcal B_{\mathcal X}\otimes\mathcal B_{\mathcal Y})$:
features and outcomes take values in the same spaces across regimes, but the
probability law governing them changes from regime to regime.
We observe $T$ data points indexed by $t=1,\dots,T$.
Regimes are indexed by $r\in\{1,\dots,R\}$ with known, non-overlapping boundaries:
regime~$r$ occupies consecutive time points
$\mathcal T_r = \{\tau_{r-1}+1,\dots,\tau_r\}$,
where $0=\tau_0 < \tau_1 < \cdots < \tau_R = T$.
Let $n_r = |\mathcal T_r| = \tau_r - \tau_{r-1}$.
Within regime~$r$, observations are i.i.d.:
$(X_t,Y_t) \stackrel{\text{i.i.d.}}{\sim} P_r$ for $t\in\mathcal T_r$,
with $P_r \neq P_s$ generally.
We write $S_r = \{(X_t,Y_t) : t\in\mathcal T_r\}$ for the sample from regime~$r$.
Throughout, we assume that the regime boundaries $\tau_0,\dots,\tau_R$ and the number of
regimes $R$ are known or are decided by the analyst. When regime membership is latent, it must first be
estimated, e.g., via structural break detection \mbox{(Bai \& Perron, 1998)} or hidden Markov
models \mbox{(Hamilton, 1989)}; we treat this as a separate, upstream inference problem and
do not address it here.

\paragraph{Sampling across regimes.}
All random variables are defined on a common probability space $(\Omega,\mathcal F,\mathbb P)$.
The samples $\{S_r\}_{r=1}^R$ are mutually independent across regimes.
Fix an index set $I\subseteq\{1,\dots,R-1\}$ corresponding to training regimes, and define the
(training) mixture distribution
\[
Q \;:=\; \sum_{r\in I} w_r P_r,
\qquad
w_r \;:=\; \frac{n_r}{\sum_{j\in I} n_j}.
\]
Let $P$ denote the deployment distribution (typically $P=P_R$).
Note that the weights $w_r$ are proportional to the regime sample sizes, so $Q$ is the empirical
training distribution that a learner actually encounters when pooling all training data. Alternative
weighting schemes (e.g., uniform $w_r=1/|I|$) define different mixture targets and may be preferable
when regime sample sizes are unbalanced or when equal importance across regimes is desired.

\paragraph{Prediction space, loss, and hypothesis classes.}
Let $(\mathcal A,\mathcal B_{\mathcal A})$ be a measurable prediction space.
Let $\ell:(\mathcal A\times\mathcal Y,\ \mathcal B_{\mathcal A}\otimes\mathcal B_{\mathcal Y})\to\R$ be measurable.
Let $\Lambda$ be a set of hyperparameters and let $\{\cF_\lambda\}_{\lambda\in\Lambda}$ be hypothesis classes with
$\cF_\lambda \subseteq \{f:(\mathcal X,\mathcal B_{\mathcal X})\to(\mathcal A,\mathcal B_{\mathcal A})\text{ measurable}\}$.


\paragraph{Existence and integrability.}
For each $\lambda\in\Lambda$, every $f\in\cF_\lambda$ is measurable $(\mathcal X,\mathcal B_{\mathcal X})\to(\mathcal A,\mathcal B_{\mathcal A})$
and has finite expected loss under the training mixture:
$\E_{(X,Y)\sim Q}\!\big[|\ell(f(X),Y)|\big]<\infty$.
Assume that for each $\lambda\in\Lambda$ the set
\[
\argmin_{f\in\cF_\lambda}\; \E_{(X,Y)\sim Q}\big[\ell(f(X),Y)\big]
\]
is nonempty; fix any element and denote it by $f^\star_{\lambda,Q}$.
That is, $f^\star_{\lambda,Q}$ is the population risk minimizer within
$\cF_\lambda$ under the training distribution~$Q$ --- the best predictor one
could identify given infinite data from~$Q$ and hyperparameter~$\lambda$.
Moreover, assume that for each $\lambda\in\Lambda$ this minimizer has finite
expected loss under the deployment distribution as well:
\[
\E_{(X,Y)\sim P}\!\big[|\ell(f^\star_{\lambda,Q}(X),Y)|\big]<\infty.
\]
Together, these assumptions guarantee that the training problem is well-posed
(a minimizer exists and the expected loss is finite) and that we can
meaningfully evaluate the trained model's risk on the deployment regime.

\paragraph{Deployment risk and oracle hyperparameter.}
Define the deployment risk
\[
\cR(\lambda;\,Q\to P)
:= \E_{(X,Y)\sim P}\big[\ell(f^\star_{\lambda,Q}(X),Y)\big],
\tag{D1}
\]
and any oracle hyperparameter
\[
\lambda^\star \in \argmin_{\lambda\in\Lambda}\ \cR(\lambda;\,Q\to P).
\tag{D2}
\]





\paragraph{Example: regularized linear regression.}
For a linear model family, each $\lambda$ corresponds to a regularization strength.
For each $\lambda$, we fit coefficients by minimizing (regularized) training loss on the
training mixture $Q$, yielding $f^\star_{\lambda,Q}$. We then select the $\lambda$ whose
resulting model has the smallest expected loss on the deployment regime $P$ -- estimated
via a regime-aware cross-validation procedure.

\paragraph{Other performance metrics.}
The formulation above uses a pointwise loss $\ell$ so that risk is an expectation under $P$.
Many evaluation criteria in practice (e.g., AUC, F1, profit/utility) are not of this form; in such cases we
either (i) replace $\ell$ by an appropriate surrogate loss whose minimizer targets the desired metric, or (ii) define a metric functional $\mathcal{M}(f;P)$ and select $\lambda$ by minimizing an empirical estimate
$\widehat{\mathcal{M}}(\hat f_{\lambda,S_{\mathrm{train}}};P)$ computed on held-out regimes using the same
regime-aware splitting scheme.



\section{Comparison to Temporal CV}

Temporal cross-validation respects the time-ordering of regimes: at each fold, we train on
past regimes and validate on a future regime. We now show that temporal CV targets a
historical average of next-block risks, which generally differs from the deployment risk
$\cR(\lambda; Q \to P)$.

\paragraph{Temporal CV procedure.}
Suppose we have $K$ folds indexed by $k = 1, \dots, K$. At fold $k$, let $I_k \subset \{1, \dots, R-1\}$
denote the set of training regimes and let $v_k \in \{2, \dots, R\}$ denote the validation regime,
where $\max(I_k) < v_k$ (training regimes precede validation). Define the fold-$k$ training set
$S^{(k)}_{\mathrm{train}} := \bigcup_{r \in I_k} S_r$ and let $\hat f_{\lambda, S^{(k)}_{\mathrm{train}}}$
denote the fitted model. The temporal CV estimator is
\[
\widehat{\mathrm{CV}}_{\mathrm{temp}}(\lambda)
\;:=\;
\frac{1}{K} \sum_{k=1}^{K}
\underbrace{\frac{1}{n_{v_k}} \sum_{t\in\mathcal T_{v_k}}
\ell\bigl(\hat f_{\lambda, S^{(k)}_{\mathrm{train}}}(X_{t}),\, Y_{t}\bigr)}_{\displaystyle =:\; \widehat{\cR}_k(\lambda)}.
\]

\paragraph{Common temporal CV schemes.}
Two standard choices for $(I_k, v_k)$ are:
\begin{itemize}
  \item \textbf{Expanding window:} $I_k = \{1, \dots, k\}$ and $v_k = k+1$. Training accumulates all
  past regimes.
  \item \textbf{Sliding window (width $w$):} $I_k = \{k, \dots, k+w-1\}$ and $v_k = k+w$. Training
  uses only the $w$ most recent regimes before validation.
\end{itemize}
Both schemes respect temporal ordering, but they induce different training distributions at each fold.
For expanding window, $k = 1, \dots, R-1$ so $K = R-1$; for sliding window of width $w$,
$k = 1, \dots, R-w$ so $K = R-w$.

\paragraph{Population target of temporal CV.}
For each fold $k$, define the fold-$k$ training mixture
\[
Q^{(k)} \;:=\; \sum_{r \in I_k} w_r^{(k)} P_r,
\qquad
w_r^{(k)} \;:=\; \frac{n_r}{\sum_{j \in I_k} n_j},
\]
and let $f^\star_{\lambda, Q^{(k)}}$ denote the population risk minimizer trained on $Q^{(k)}$.
The population target of temporal CV is
\[
\overline{\mathrm{CV}}_{\mathrm{temp}}(\lambda)
\;:=\;
\frac{1}{K} \sum_{k=1}^{K} \cR(\lambda;\, Q^{(k)} \to P_{v_k}),
\]
where $\cR(\lambda; Q^{(k)} \to P_{v_k}) := \E_{(X,Y) \sim P_{v_k}}[\ell(f^\star_{\lambda, Q^{(k)}}(X), Y)]$.
By standard arguments (the law of large numbers applied to validation losses, with training and
validation regimes independent), temporal CV converges to this target as sample sizes grow.

\paragraph{When does the mismatch matter?}
The gap between $\overline{\mathrm{CV}}_{\mathrm{temp}}(\lambda)$ and
$\cR(\lambda;\,Q\to P)$ depends on the dynamics of the regime process.
If regimes are drawn i.i.d.\ from a stationary distribution, then
the validation distributions $P_{v_k}$ are exchangeable and temporal CV
targets an average risk that is already a natural pre-deployment objective;
in this case the mismatch is benign.
The distinction becomes consequential when regimes exhibit
\emph{persistence}: if the current regime is informative about the next,
the deployment distribution $P$ is partially predictable from recent
regime history, and a CV procedure that ignores this structure
systematically mis-weights its validation evidence.
The regime-aware procedure developed below is therefore most relevant in
the persistent-regime setting.


\section{Regime-Aware Cross-Validation}



\end{document}
